%%% variant: default
\subsection{Wärmeleitung}

Wie hängt der Wärmestrom von der Wanddicke ab?

\par
Bei gleichem Temperaturunterschied halbiert sich der Wärmestrom, wenn sich die Wanddicke verdoppelt.
\par

\par
\par\begingroup\centering\noindent\begin{adjustbox}{width=12em,max width=\linewidth}

\begin{tikzpicture}
  \draw[->] (0,0) -- (3,0) node[right] {$x$};
  \draw[->] (0,0) -- (0,2) node[above] {$T$};
  \draw[thick] (0,1.7) -- (2.8,0.3);
\end{tikzpicture}

\end{adjustbox}\par
Temperaturverlauf in einer ebenen Wand bei stationärer Wärmeleitung.
\par\endgroup
\par

Die Wärmeleitfähigkeit sei konstant.

%%% variant: instructor
\subsection{Wärmeleitung}

Wie hängt der Wärmestrom von der Wanddicke ab?

\par
Bei gleichem Temperaturunterschied halbiert sich der Wärmestrom, wenn sich die Wanddicke verdoppelt.
\par

\par
\par\begingroup\centering\noindent\begin{adjustbox}{width=12em,max width=\linewidth}

\begin{tikzpicture}
  \draw[->] (0,0) -- (3,0) node[right] {$x$};
  \draw[->] (0,0) -- (0,2) node[above] {$T$};
  \draw[thick] (0,1.7) -- (2.8,0.3);
\end{tikzpicture}

\end{adjustbox}\par
Temperaturverlauf in einer ebenen Wand bei stationärer Wärmeleitung.
\par\endgroup
\par

Die Wärmeleitfähigkeit sei konstant.

%%% variant: handout
\subsection{Wärmeleitung}

Wie hängt der Wärmestrom von der Wanddicke ab?

\par
Bei gleichem Temperaturunterschied halbiert sich der Wärmestrom, wenn sich die Wanddicke verdoppelt.
\par

\par
\par\begingroup\centering\noindent\begin{adjustbox}{width=12em,max width=\linewidth}

\begin{tikzpicture}
  \draw[->] (0,0) -- (3,0) node[right] {$x$};
  \draw[->] (0,0) -- (0,2) node[above] {$T$};
  \draw[thick] (0,1.7) -- (2.8,0.3);
\end{tikzpicture}

\end{adjustbox}\par
Temperaturverlauf in einer ebenen Wand bei stationärer Wärmeleitung.
\par\endgroup
\par

Die Wärmeleitfähigkeit sei konstant.

%%% variant: slides
\begin{frame}[t]
\frametitle{Wärmeleitung}

Wie hängt der Wärmestrom von der Wanddicke ab?

\begin{columns}[T,onlytextwidth]
\begin{column}{\dimexpr(\linewidth-1\columnsep)/2\relax}
Bei gleichem Temperaturunterschied halbiert sich der Wärmestrom, wenn sich die Wanddicke verdoppelt.
\end{column}

\begin{column}{\dimexpr(\linewidth-1\columnsep)/2\relax}
\par\begingroup\centering\noindent\begin{adjustbox}{width=12em,max width=\linewidth}

\begin{tikzpicture}
  \draw[->] (0,0) -- (3,0) node[right] {$x$};
  \draw[->] (0,0) -- (0,2) node[above] {$T$};
  \draw[thick] (0,1.7) -- (2.8,0.3);
\end{tikzpicture}

\end{adjustbox}\par
\par\endgroup
\end{column}
\end{columns}

Die Wärmeleitfähigkeit sei konstant.

\end{frame}

%%% variant: slides-dozent-h1-frame
\begin{frame}[t]
\frametitle{Wärmeleitung}

Wie hängt der Wärmestrom von der Wanddicke ab?

\begin{columns}[T,onlytextwidth]
\begin{column}{\dimexpr(\linewidth-1\columnsep)/2\relax}
Bei gleichem Temperaturunterschied halbiert sich der Wärmestrom, wenn sich die Wanddicke verdoppelt.
\end{column}

\begin{column}{\dimexpr(\linewidth-1\columnsep)/2\relax}
\par\begingroup\centering\noindent\begin{adjustbox}{width=12em,max width=\linewidth}

\begin{tikzpicture}
  \draw[->] (0,0) -- (3,0) node[right] {$x$};
  \draw[->] (0,0) -- (0,2) node[above] {$T$};
  \draw[thick] (0,1.7) -- (2.8,0.3);
\end{tikzpicture}

\end{adjustbox}\par
\par\endgroup
\end{column}
\end{columns}

Die Wärmeleitfähigkeit sei konstant.

\end{frame}

%%% variant: slides-h1-hide-level-1
\begin{frame}[t]
\par\medskip\noindent\textbf{Wärmeleitung}\par


Wie hängt der Wärmestrom von der Wanddicke ab?

\begin{columns}[T,onlytextwidth]
\begin{column}{\dimexpr(\linewidth-1\columnsep)/2\relax}
Bei gleichem Temperaturunterschied halbiert sich der Wärmestrom, wenn sich die Wanddicke verdoppelt.
\end{column}

\begin{column}{\dimexpr(\linewidth-1\columnsep)/2\relax}
\par\begingroup\centering\noindent\begin{adjustbox}{width=12em,max width=\linewidth}

\begin{tikzpicture}
  \draw[->] (0,0) -- (3,0) node[right] {$x$};
  \draw[->] (0,0) -- (0,2) node[above] {$T$};
  \draw[thick] (0,1.7) -- (2.8,0.3);
\end{tikzpicture}

\end{adjustbox}\par
\par\endgroup
\end{column}
\end{columns}

Die Wärmeleitfähigkeit sei konstant.

\end{frame}

